\documentclass[11pt]{article}
\usepackage[margin=1in]{geometry}
\usepackage{booktabs}
\usepackage{array}
\usepackage[T1]{fontenc}
\usepackage[utf8]{inputenc}
\usepackage{lmodern}
\setlength{\parskip}{0.6em}
\setlength{\parindent}{0pt}
\title{Harborlight Savings CECL Gap Assessment}
\date{March 13, 2026}
\begin{document}
\maketitle
\section*{Executive Summary}
This assessment covers Harborlight Savings's Q1 2026 CECL Readiness Gap Assessment package. The materials support a meaningful documentation-led review, but the absence of a runnable reserve engine and lineaged execution evidence prevents an execution-based opinion. The review therefore focuses on package sufficiency, scenario consistency, segment reconciliation, and qualitative overlay support.
\section*{Package Reviewed}
\begin{itemize}
\item Methodology, model overview, scenario assumptions, overlay memo, and prior review note
\item Scenario tables and provided reserve outputs
\item Evidence request log and governance/gap-tracker materials
\end{itemize}
\section*{Work Performed}
\begin{itemize}
\item Inventoried the supplied package and identified execution blockers
\item Compared scenario narrative to supplied scenario tables
\item Compared documented segment structure to the provided reserve outputs
\item Assessed whether overlay documentation was consistent with the supplied reserve bridge
\end{itemize}
\section*{Provided Reserve Outputs}
\begin{tabular}{p{0.24\textwidth} | p{0.24\textwidth} | p{0.24\textwidth}}
\toprule
Scenario & Reserve & Reserve Rate \\
\midrule
Baseline & \$38,420,000 & 1.84\% \\
Adverse & \$52,180,000 & 2.48\% \\
Severe & \$53,010,000 & 2.51\% \\
\bottomrule
\end{tabular}
\section*{Key Gaps and Findings}
\begin{itemize}
\item [HIGH] Execution-based review is blocked by missing reserve engine and lineage evidence: The package includes prior outputs and narrative materials, but no reserve engine, reproducibility script, or lineaged runbook sufficient for execution-based review.
\item [HIGH] Scenario narrative is not fully aligned to the supplied scenario tables: The severe scenario narrative describes a uniformly harsher path, but the supplied numeric scenarios show 4 quarter(s) where severe house-price growth is less severe than adverse.
\item [MEDIUM] Documented segment structure does not reconcile to the provided reserve outputs: Documentation describes 5 segments, while the provided reserve outputs reconcile to 4 output segments.
\item [MEDIUM] Overlay documentation understates the magnitude implied by the provided reserve bridge: The overlay memo frames qualitative adjustment as capped at 6.0 bps, but the provided bridge shows up to 18.0 bps.
\end{itemize}
\section*{Coverage Position}
This package supports a documentation-led gap assessment only. Baseline reproduction, scenario reruns, sensitivity testing, and implementation validation remain blocked pending runnable model artifacts and execution lineage.
\section*{Evidence Requests}
\begin{itemize}
\item Provide the runnable CECL reserve engine or reproducibility notebook used to generate the supplied reserve outputs.
\item Provide a lineaged runbook linking the prior output snapshots to the model version and scenario definitions reviewed by governance.
\item Reconcile documented segment taxonomy to the segment structure used in the reserve outputs.
\item Refresh scenario documentation so the severe narrative matches the numeric severe path quarter by quarter.
\end{itemize}
\section*{Overall Assessment}
The package is sufficient for a structured gap assessment but not for a model-driven CECL review opinion. The primary blockers are missing implementation artifacts, unresolved scenario-definition inconsistencies, and overlay support that is not quantitatively reconciled.
\end{document}